% =========================================================
% Distance Between Sets
% Source: volume-iii/analysis/metric-spaces/notes/metric-spaces/
%
% Dependencies:
%   notes-metric-space.tex      — def:metric, def:open-ball
%   notes-properties-metric-space.tex — def:diameter, def:bounded-set
%   notes-limit-points.tex      — def:closure (for the closure characterization)
% =========================================================

% ---------------------------------------------------------
\subsubsection{Distance Between Sets}
\label{sec:distance-between-sets}
% ---------------------------------------------------------

% ---- Toolkit ----
\begin{tcolorbox}[colback=gray!6, colframe=gray!40, arc=2pt,
  left=6pt, right=6pt, top=4pt, bottom=4pt,
  title={\small\textbf{Distance Between Sets --- Quick Reference}},
  fonttitle=\small\bfseries]
\small
\begin{tabular}{l l l l}
\toprule
\textbf{Label} & \textbf{Statement} & \textbf{Proof method} & \textbf{Detail} \\
\midrule
D\ref{def:dist-point-set}
  & $d(x, A) := \inf\{d(x,a) : a \in A\}$
  & Definition
  & \hyperref[def:dist-point-set]{$\downarrow$} \\[4pt]
D\ref{def:dist-sets}
  & $d(A, B) := \inf\{d(a,b) : a \in A,\, b \in B\}$
  & Definition
  & \hyperref[def:dist-sets]{$\downarrow$} \\[4pt]
P\ref{prop:dist-point-set-closure}
  & $d(x, A) = 0 \iff x \in \overline{A}$
  & Sequential char.\ of closure
  & \hyperref[prop:dist-point-set-closure]{$\downarrow$} \\[4pt]
P\ref{prop:dist-sets-not-metric}
  & $d(A,B) = 0 \not\Rightarrow A \cap B \neq \varnothing$
  & Counterexample
  & \hyperref[prop:dist-sets-not-metric]{$\downarrow$} \\[4pt]
P\ref{prop:dist-sets-symmetric}
  & $d(A, B) = d(B, A)$
  & Symmetry of $d$
  & \hyperref[prop:dist-sets-symmetric]{$\downarrow$} \\[4pt]
\bottomrule
\end{tabular}
\end{tcolorbox}

% ---------------------------------------------------------

The diameter of a set measures its internal spread --- the largest distance
\emph{within} a set.
We now introduce a companion notion: the distance \emph{between} two sets,
which measures how far apart they are from each other.

% ---- Definition: Distance from a Point to a Set ----
\begin{tcolorbox}[colback=propbox, colframe=propborder, arc=2pt,
  left=6pt, right=6pt, top=4pt, bottom=4pt,
  title={\small\textbf{Definition (Distance from a Point to a Set)}},
  fonttitle=\small\bfseries]
\begin{definition}[Distance from a Point to a Set]
\label{def:dist-point-set}
Let $(X, d)$ be a metric space, let $A \subseteq X$ be non-empty,
and let $x \in X$.
The \emph{distance from $x$ to $A$} is
\[
d(x, A) \;:=\; \inf\{\, d(x, a) : a \in A \,\}.
\]
\end{definition}
\end{tcolorbox}

\begin{remark*}[Fully quantified form]
\[
\forall\,(X,d)\;\text{metric space},\;
\forall\,A \subseteq X\;\text{non-empty},\;
\forall\,x \in X,\quad
d(x,A) := \inf\{\, d(x,a) : a \in A \,\}.
\]
\end{remark*}

\begin{remark*}[Intuition]
$d(x, A)$ is the closest $A$ gets to $x$.
It is the infimum --- not necessarily a minimum --- of all distances from $x$
to points of $A$.
The infimum may not be achieved: if $x \notin A$ and $A$ is open, there may
be no $a \in A$ realising the infimum.
In $\mathbb{R}$, for instance, $d(0, (0,1)) = 0$ but no point of $(0,1)$
is at distance $0$ from $0$.
\end{remark*}

% ---- Definition: Distance Between Sets ----
\begin{tcolorbox}[colback=propbox, colframe=propborder, arc=2pt,
  left=6pt, right=6pt, top=4pt, bottom=4pt,
  title={\small\textbf{Definition (Distance Between Sets)}},
  fonttitle=\small\bfseries]
\begin{definition}[Distance Between Sets]
\label{def:dist-sets}
Let $(X,d)$ be a metric space and let $A, B \subseteq X$ be non-empty.
The \emph{distance between $A$ and $B$} is
\[
d(A, B) \;:=\; \inf\{\, d(a, b) : a \in A,\; b \in B \,\}.
\]
\end{definition}
\end{tcolorbox}

\begin{remark*}[Fully quantified form]
\[
\forall\,(X,d)\;\text{metric space},\;
\forall\,A,B \subseteq X\;\text{non-empty},\quad
d(A,B) := \inf\{\, d(a,b) : a \in A,\; b \in B \,\}.
\]
\end{remark*}

\begin{remark*}[Relation to point-to-set distance]
The two definitions are related by
\[
d(A, B) \;=\; \inf_{a \in A}\, d(a, B) \;=\; \inf_{b \in B}\, d(b, A).
\]
The distance between sets is the infimum of all distances from points of
one set to the other set.
\end{remark*}

% ---- Proposition: d(x,A) = 0 iff x in closure(A) ----

\begin{proposition}[Distance zero to a set characterizes closure]
\label{prop:dist-point-set-closure}
Let $(X,d)$ be a metric space and let $A \subseteq X$ be non-empty.
For any $x \in X$:
\[
d(x, A) = 0
\;\iff\;
x \in \overline{A}.
\]
\end{proposition}

\begin{proof}
Unfolding definitions:
\begin{align*}
d(x, A) = 0
&\iff \inf\{\, d(x,a) : a \in A \,\} = 0 \\
&\iff \forall\, r > 0,\; \exists\, a \in A \text{ with } d(x,a) < r \\
&\iff \forall\, r > 0,\; B_r(x) \cap A \neq \varnothing \\
&\iff x \in \overline{A},
\end{align*}
where the last equivalence is the geometric characterization of closure
(Theorem~\ref{thm:closure-equivalences}).
\end{proof}

\begin{remark*}[Consequence: a fourth characterization of closure]
Proposition~\ref{prop:dist-point-set-closure} adds a fourth language to
the unifying closure theorem:
\[
x \in \overline{A}
\;\iff\;
\forall r>0,\; B_r(x) \cap A \neq \varnothing
\;\iff\;
\exists (a_n) \subseteq A,\; a_n \to x
\;\iff\;
d(x, A) = 0.
\]
The metric-valued characterization $d(x,A) = 0$ is often the most
computationally convenient: checking $d(x,A) = 0$ reduces to an
infimum calculation rather than an existential statement about sequences.
\end{remark*}

\begin{remark*}[Common error]
$d(x, A) = 0$ does \emph{not} imply $x \in A$ unless $A$ is closed.
The infimum of distances may equal zero because $x$ can be approximated
arbitrarily closely by points of $A$, even if $A$ does not contain $x$ itself.
This is precisely the distinction between $A$ and $\overline{A}$.
\end{remark*}

% ---- Proposition: d(A,B) = 0 does not imply A ∩ B ≠ ∅ ----

\begin{proposition}[$d(A,B) = 0$ does not imply intersection]
\label{prop:dist-sets-not-metric}
In general, $d(A, B) = 0$ does not imply $A \cap B \neq \varnothing$.
\end{proposition}

\begin{example}[Disjoint sets at distance zero]
\label{ex:dist-sets-zero-disjoint}
In $(\mathbb{R}, |\cdot|)$, let $A = (0,1)$ and $B = (1,2)$.
These sets are disjoint: $A \cap B = \varnothing$.
Yet
\[
d(A, B) = \inf\{\, |a - b| : a \in (0,1),\; b \in (1,2) \,\} = 0,
\]
since the infimum is not achieved --- points in $A$ can get arbitrarily
close to $1$, and points in $B$ can get arbitrarily close to $1$,
but neither set contains $1$.
\end{example}

\begin{remark*}[Why $d(A,B)$ is not a metric on sets]
The failure $d(A,B) = 0 \not\Rightarrow A = B$ (or even $A \cap B \neq \varnothing$)
means that $d$ does not satisfy the identity axiom for a metric.
The function $(A,B) \mapsto d(A,B)$ is therefore \emph{not} a metric on
the collection of subsets of $X$.

The correct replacement --- a genuine metric on the collection of non-empty
compact subsets of a metric space --- is the \emph{Hausdorff metric}:
\[
d_H(A, B)
:= \max\Bigl(
  \sup_{a \in A} d(a, B),\;
  \sup_{b \in B} d(b, A)
\Bigr).
\]
The Hausdorff metric is a topic for later study (functional analysis or
advanced topology).
\end{remark*}

% ---- Proposition: symmetry ----

\begin{proposition}[Symmetry of set distance]
\label{prop:dist-sets-symmetric}
For non-empty $A, B \subseteq X$:
\[
d(A, B) = d(B, A).
\]
\end{proposition}

\begin{proof}
$d(A,B) = \inf\{d(a,b) : a \in A, b \in B\}
= \inf\{d(b,a) : b \in B, a \in A\}
= d(B,A)$,
using symmetry of $d$ (axiom M3).
\end{proof}

% ---- Worked Example ----

\begin{example}[Computing $d(x, A)$ and $d(A, B)$ in $\mathbb{R}$]
\label{ex:dist-point-set-concrete}
Work in $(\mathbb{R}, |\cdot|)$.

\medskip
\noindent\textbf{(a) $d(3, [0,1])$.}

$d(3, [0,1]) = \inf\{|3 - a| : a \in [0,1]\}$.
Since $3 > 1$, the infimum is achieved at $a = 1$:
$d(3, [0,1]) = |3 - 1| = 2$.

\medskip
\noindent\textbf{(b) $d(0.5, [0,1])$.}

$0.5 \in [0,1]$, so $d(0.5, [0,1]) = 0$.
(The infimum is achieved at $a = 0.5$ itself.)

\medskip
\noindent\textbf{(c) $d([0,1], [3,4])$.}

$d([0,1],[3,4]) = \inf\{|a - b| : a \in [0,1], b \in [3,4]\}$.
The two sets are on opposite sides of the gap $(1,3)$.
The infimum is achieved at $a = 1$, $b = 3$:
$d([0,1],[3,4]) = |1 - 3| = 2$.

\medskip
\noindent\textbf{(d) $d((0,1), (1,2))$.}

Both sets are open, sharing the boundary point $1$ that belongs to neither.
$d((0,1),(1,2)) = 0$ (infimum not achieved; see Example~\ref{ex:dist-sets-zero-disjoint}).
\end{example}

\begin{remark*}[When is the infimum achieved?]
$d(x, A)$ is achieved (i.e.\ is a minimum) when $A$ is closed and bounded
in a complete metric space --- more precisely, when $A$ is compact.
Similarly, $d(A, B)$ is achieved when both $A$ and $B$ are compact and disjoint.
These results belong to the compactness chapter.
For now, treat the infimum as genuinely possibly unachieved.
\end{remark*}
